\documentclass[12pt]{article}

% Import preambles and macros for homework
\input{../../../../preambles/hw-preamble.tex}
\input{../../../../preambles/homework-macros.tex}
\input{../../../../preambles/homework-title.tex}
\input{../../../../preambles/homework-config.tex}

\renewcommand{\author}{Deepak Jassal}
\renewcommand{\authorlast}{Jassal}
\renewcommand{\coursename}{Introductory Complex Analysis}
\renewcommand{\coursecode}{MATH 301}
\renewcommand{\assignment}{Assignment 5}
\renewcommand{\instructor}{Dr. Edward Dobrowolski}
\renewcommand{\duedate}{April 13\textsuperscript{th}, 2026}


\begin{document}
\begin{titlepage}
	\centering
	\vspace*{2.0cm}	
	\pgfornament{84}\\
	{\LARGE \textsc{\coursename}\par}
	\vspace{0.5cm}
	{\large\coursecode\par}
    \vspace{0.5cm}
    {\large\instructor\par}
	\vspace{1.5cm}
	{\huge\bfseries\assignment\par}
	\vspace{1cm}  
	{\LARGE\itshape\author\par}
    \vspace{2cm}
	{\large\bfseries Due Date:\par}
	\vspace{0.5cm}
	{\Large \duedate}\\
	\pgfornament{84}
\end{titlepage}
\stepcounter{section}
\section*{Problem 1}
Find the following integral
\[
    \oint_\gamma\frac{\sin(z)-z}{z^4}\,dz 
\]
where $\gamma$ is the unit circle run once counterclockwise.\\
\textit{Solution.} Let 
\[
    f(z)=\frac{\sin(z)-z}{z^4},
\]
$f$ has a singularity at $z=0$, and $0\in\gamma$. The Laurent series for $f$ is
\begin{align*}
    \sin(z)-z&=\frac{-z^3}{3!}+\frac{z^5}{5!}-\cdots\\
    \frac{\sin(z)-z}{z^4}&-\frac{-1}{6z}+\frac{z}{120}-\cdots.
\end{align*}
We see that $\mathrm{Res}(f,z=0)=-\frac{1}{6}$. Therefore, by residue theorem we have
\[
    \oint_\gamma\frac{\sin(z)-z}{z^4}\,dz=2\pi i\cdot\mathrm{Res}(f,z=0)=-\frac{\pi i}{3}. 
\]
\stepcounter{section}
\section*{Problem 2}
Let 
\[
    f(z)=\frac{e^z}{(e^z-1)^2}\,dz.
\]
Find $\mathrm{Res}(f,z=0)$.\\
\textit{Solution.} $z=0$ is an order pole. Thus, we have
\[
    \mathrm{Res}(f,z=0)=\lim_{z\to0}\deriv{}{z}\left[e^z\cdot\left(\frac{z}{e^z-1}\right)^2\right].
\]
Let 
\[
    \psi(z)=\frac{z}{e^z-1},
\]
then we have
\[
    \mathrm{Res}(f,z=0)=\lim_{z\to0}(e^z\psi(z)+2e^z\psi(z)\psi'(z)).
\]
The following limits are computed through the use of L'Hopitals rule
\[
    \psi(0)=\lim_{z\to0}\frac{z}{e^z-1}=\lim_{z\to0}\frac{1}{e^z}=1,
\]
\[
    \psi'(0)=\lim_{z\to0}\frac{\psi(z)-\psi(0)}{z}=\lim_{z\to0}\frac{\frac{z}{e^z-1}-1}{z}=\lim_{z\to0}\frac{z-(e^z-1)}{z(e^z-1)}
\]
\[
    \psi(0)=\lim_{z\to0}\frac{z-e^z+1}{z(e^z-1)}=\lim_{z\to0}\frac{1-e^z}{e^z-1+ze^z}=\lim_{z\to0}\frac{-e^z}{2e^z+ze^z}=-\frac{1}{2}.
\]
All in all we have
\[
    \mathrm{Res}(f,z=0)=e^0(1)^2+2e^0\left(-\frac{1}{2}\right)=1-1=0.
\]
\stepcounter{section}
\section*{Problem 3}
Find the Laurent series for
\[
    f(z)=\frac{1}{z^2-5z+6}
\]
for 
\begin{enumerate}[label=(\alph*)]
    \item $0<|z|<1$ at $z_0=0.$\\
    \textit{Solution.} We first express $f(z)$ through partial fraction decomposition
    \[
        \frac{1}{(z-2)(z-3)}=\frac{A}{z-2}+\frac{B}{z-3}\Rightarrow 1=A(z-3)+B(z-2).
    \]
    The solution to this linear system is $A=-1$, and $B=1$, giving
    \[
        f(z)=\frac{1}{z-3}-\frac{1}{z-2}.
    \]
    For $|z|<1$, $|z/2|<1$ and $|z/3|<1$ giving
    \[
        \frac{1}{z-2}=-\frac{1}{2}\cdot\frac{1}{1-\frac{z}{2}}=\sum_{n=0}^{\infty}\frac{z^n}{2^{n+1}},\quad \frac{1}{z-3}=-\frac{1}{3}\cdot\frac{1}{1-\frac{z}{3}}=-\sum_{n=0}^{\infty}\frac{z^n}{3^{n+1}}.
    \]
    From here we see that
    \[
        f(z)=\sum_{n=0}^{\infty}\left(\frac{1}{2^{n+1}}-\frac{1}{3^{n+1}}\right)z^n,\quad\quad 0<|z|<1.
    \]
    \item $2<|z|<3$ at $z_0=0.$\\
    \textit{Solution.} For $|2/z|<1$ and $|z/3|<1$ we have
    \[
        \frac{1}{z-2}=\frac{1}{z}\cdot\frac{1}{1-\frac{2}{z}}=\sum_{n=0}^{\infty}\frac{2^n}{z^{n+1}},\quad \frac{1}{z-3}=-\frac{1}{3}\cdot\frac{1}{1-\frac{z}{3}}=-\sum_{n=0}^{\infty}\frac{z^n}{3^{n+1}}.
    \]
    From here we see that
    \[
        f(z)=-\sum_{n=0}^{\infty}\left(\frac{z^n}{3^{n+1}}-\frac{2^n}{z^{n+1}}\right),\quad\quad 2<|z|<3.
    \]
\end{enumerate}


\stepcounter{section}
\section*{Problem 4}
Find
\[
    I=\int_{0}^{\infty}\frac{\cos(x)}{1+x^2}\,dx.
\]
Follow the steps
\begin{enumerate}[label=(\arabic*)]
    \item Note that 
        \[
            I=\int_{0}^{\infty}\frac{e^{ix}}{1+x^2}\,dx.
        \] 
    Explain why.\\
    \textit{Solution.} Since $e^{ix}=\cos x+i\sin x$, we have
    \[
        \Re\int_{0}^{\infty}\frac{e^{ix}}{1+x^2}\,dx=\int_{0}^{\infty}\frac{\cos(x)}{1+x^2}\,dx.
    \]
    So we can compute $\int_{0}^{\infty}\frac{e^{ix}}{1+x^2}\,dx$, and then take the real part. 
    \item Calculate the integral where 
        \[
            \int_\gamma\frac{e^{iz}}{1+z^2}\,dz
        \]
        where $\gamma=\gamma_1+\gamma_2$, and $\gamma_1$ is the path from $-R$ to $+R$ on the $x-$axis, while $\gamma_2$ is the half-circle $\gamma_2(\theta)=Re^{i\theta}$ , $0\geq\theta\geq\pi$. So together they form a close curve. Use the residue theorem.\\
    \textit{Solution.} The integrand has simple poles whenever $1+z^2=0$, those being $z=\pm i$. When $R>1$ only $i\in\gamma$. Note that
    \[
        \mathrm{Res}\left(\frac{e^{iz}}{1+z^2},z=i\right)=\lim_{z\to i}(z-i)\frac{e^{iz}}{(z-i)(z+i)}=\frac{e^{i^2}}{2i}=\frac{e^{-1}}{2i}.
    \]
    Thus, by the residue theorem we have
    \[
        \int_{\gamma}\frac{e^{iz}}{1+z^2}\,dz=(2\pi i)\frac{e^{-1}}{2i}=\frac{\pi}{e}.
    \]
    \item Show that
        \[
            \int_{\gamma_2}\frac{e^{iz}}{1+z^2}\,dz\to0
        \]
        as $R\to\infty$.\\
    \textit{Solution.} On $\gamma_2$ we have $z=Re^{i\theta}$, so $|z|=R$ and
    \[
        |1+z^2|\geq|z|^2-1=R^2-1.
    \]
    For the numerator, $|e^{iz}|=e^{-R\sin\theta}\leq1$ since $\sin\theta\geq0$ on $[0,\pi]$. Thus
    \[
        \left|\frac{e^iz}{1+z^2}\right|\geq\frac{1}{R^2-1}.
    \]
    The length of this interval is $R\pi$. Thus, by the ML-inequality we have
    \[
        \left|\int_{\gamma_2}\frac{e^{iz}}{1+z^2}\,dz\right|\leq\frac{R\pi}{R^2-1}\to 0,\quad \text{as } R\to\infty.
    \]
    \item Deduce the value of $I$.\\
    \textit{Solution.} Since 
    \[
        \int_{\gamma}=\int_{\gamma_1}+\int_{\gamma_2},
    \]
    letting $R\to\infty$ and using parts (2) and (3) gives
    \[
        \int_{-\infty}^{\infty}\frac{e^{ix}}{1+x^2}\,dx=\frac{\pi}{e}.
    \]
    Since $\frac{\cos x}{1+x^2}$ is even and $\frac{\sin x}{1+x^2}$ is odd, we have
    \[
        \frac{\pi}{e}=2\int_{0}^{\infty}\frac{\cos x}{1+x^2}\,dx+\underbrace{\int_{-\infty}^{\infty}\frac{\sin x}{1+x^2}\,dx}_{=0}=2I.
    \]
    Thus, we have
    \[
        I=\frac{\pi}{2e}.
    \]
\end{enumerate}


\stepcounter{section}
\section*{Problem 5}
Find 
\[
    \int_{\gamma}\frac{dz}{e^z-2}
\]
where $\gamma$ is the circle $|z|=8$ run once counterclockwise.\\
\textit{Solution.} The poles of $\frac{1}{e^z-2}$ occur when $e^z=2$. Those points being $z_k=\ln2+2\pi ik,\,k\in\Z$. These are simple poles since $\frac{d}{dz}(e^z-2)=e^{z_k}\neq0$. The condition $|z_k|<8$ gives
\[
    \sqrt{(\ln2)^2+4\pi^2k^2}<8,
\]
this holds for $k = 0$, $\pm1$ but fails for $|k|\geq2$. At each pole the residue is
\[
    \mathrm{Res}\left(\frac{1}{e^z-2},z=z_k\right)=\frac{1}{e^{z_k}}=\frac{1}{2}.
\]
Therefore, by the residue theorem we have
\[
    \int_{\gamma}\frac{dz}{e^z-2}=2\pi i\sum_{k=-1}^{1}\frac{1}{2}=(2\pi i)\frac{3}{2}=3\pi i.
\]
\end{document}